% OpenStax Calculus Volume 2 -- Section 3.2 Exercises: Trigonometric Integrals
% Exercises reproduced from OpenStax, licensed under CC BY 4.0.
% Source: https://openstax.org/books/calculus-volume-2/pages/3-2-trigonometric-integrals
\documentclass{exercises}
\title{OpenStax Calculus Volume 2}
\subtitle{Section 3.2 Exercises: Trigonometric Integrals}
\sourceurl{https://openstax.org/books/calculus-volume-2/pages/3-2-trigonometric-integrals}
\author{}
\date{}

\begin{document}
\maketitle


Fill in the blank to make a true statement.

\begin{enumerate}[start=1]
\item $\sin^{2}x+\underline{\hspace{1.5cm}}=1$
\item $\sec^{2}x-1=\underline{\hspace{1.5cm}}$
\end{enumerate}

Use an identity to reduce the power of the trigonometric function to a trigonometric function raised to the first power.

\begin{enumerate}[resume]
\item $\sin^{2}x=\underline{\hspace{1.5cm}}$
\item $\cos^{2}x=\underline{\hspace{1.5cm}}$
\end{enumerate}

Evaluate each of the following integrals by u-substitution.

\begin{enumerate}[resume]
\item $\int \sin^{3}x\cos x\,dx$
\item $\int \sqrt{\cos x}\sin x\,dx$
\item $\int \tan^{5}(2x)\sec^{2}(2x)\,dx$
\item $\int \sin^{7}(2x)\cos(2x)\,dx$
\item $\int \tan(\frac{x}{2})\sec^{2}(\frac{x}{2})\,dx$
\item $\int \tan^{2}x\sec^{2}x\,dx$
\end{enumerate}

Compute the following integrals using the guidelines for integrating powers of trigonometric functions. Use a CAS to check the solutions. (Note: Some of the problems may be done using techniques of integration learned previously.)

\begin{enumerate}[resume]
\item $\int \sin^{3}x\,dx$
\item $\int \cos^{3}x\,dx$
\item $\int \sin x\cos x\,dx$
\item $\int \cos^{5}x\,dx$
\item $\int \sin^{5}x\cos^{2}x\,dx$
\item $\int \sin^{3}x\cos^{3}x\,dx$
\item $\int \sqrt{\sin x}\cos x\,dx$
\item $\int \sqrt{\sin x}\cos^{3}x\,dx$
\item $\int \sec x\tan x\,dx$
\item $\int \tan(5x)\,dx$
\item $\int \tan^{2}x\sec x\,dx$
\item $\int \tan x\sec^{3}x\,dx$
\item $\int \sec^{4}x\,dx$
\item $\int \cot x\,dx$
\item $\int \csc x\,dx$
\item $\int \frac{\tan^{3}x}{\sqrt{\sec x}}\,dx$
\end{enumerate}

For the following exercises, find a general formula for the integrals.

\begin{enumerate}[resume]
\item $\int \sin^{2}(ax)\cos(ax)\,dx$
\item $\int \sin(ax)\cos(ax)\,dx$.
\end{enumerate}

Use the double-angle formulas to evaluate the following integrals.

\begin{enumerate}[resume]
\item $\int_{0}^{\pi}\sin^{2}x\,dx$
\item $\int_{0}^{\pi}\sin^{4}x\,dx$
\item $\int \cos^{2}3x\,dx$
\item $\int \sin^{2}x\cos^{2}x\,dx$
\item $\int \sin^{2}x\,dx+\int \cos^{2}x\,dx$
\item $\int \sin^{2}x\cos^{2}(2x)\,dx$
\end{enumerate}

For the following exercises, evaluate the definite integrals. Express answers in exact form whenever possible.

\begin{enumerate}[resume]
\item $\int_{0}^{2\pi}\cos x\sin 2x\,dx$
\item $\int_{0}^{\pi}\sin 3x\sin 5x\,dx$
\item $\int_{0}^{\pi}\cos(99x)\sin(101x)\,dx$
\item $\int_{-\pi}^{\pi}\cos^{2}(3x)\,dx$
\item $\int_{0}^{2\pi}\sin x\sin(2x)\sin(3x)\,dx$
\item $\int_{0}^{4\pi}\cos(x/2)\sin(x/2)\,dx$
\item $\int_{\pi/6}^{\pi/3}\frac{\cos^{3}x}{\sqrt{\sin x}}\,dx$ (Round this answer to three decimal places.)
\item $\int_{-\pi/3}^{\pi/3}\sqrt{\sec^{2}x-1}\,dx$
\item $\int_{0}^{\pi/2}\sqrt{1-\cos(2x)}\,dx$
\item Find the area of the region bounded by the graphs of the equations $y=\sin x$, $y=\sin^{3}x$, $x=0$, and $x=\frac{\pi}{2}$.
\item Find the area of the region bounded by the graphs of the equations $y=\cos^{2}x$, $y=\sin^{2}x$, $x=-\frac{\pi}{4}$, and $x=\frac{\pi}{4}$.
\item A particle moves in a straight line with the velocity function $v(t)=\sin(\omega t)\cos^{2}(\omega t)$. Find its position function $x=f(t)$ if $f(0)=0$.
\item Find the average value of the function $f(x)=\sin^{2}x\cos^{3}x$ over the interval $[-\pi,\pi]$.
\end{enumerate}

For the following exercises, solve the differential equations.

\begin{enumerate}[resume]
\item $\frac{dy}{dx}=\sin^{2}x$. The curve passes through point $(0,0)$.
\item $\frac{dy}{d\theta}=\sin^{4}(\pi \theta)$
\item Find the length of the curve $y=\ln(\csc x),\frac{\pi}{4}\le x\le \frac{\pi}{2}$.
\item Find the length of the curve $y=\ln(\sin x),\frac{\pi}{3}\le x\le \frac{\pi}{2}$.
\item Find the volume generated by revolving the curve $y=\cos(3x)$ about the $x$-axis, $0\le x\le \frac{\pi}{36}$.
\end{enumerate}

For the following exercises, use this information: The inner product of two functions $f$ and $g$ over $[a,b]$ is defined by $f(x)\cdot g(x)=\langle f,g\rangle=\int_{a}^{b}f\cdot g\,dx$. Two distinct functions $f$ and $g$ are said to be orthogonal if $\langle f,g\rangle=0$.

\begin{enumerate}[resume]
\item Show that $\{\sin(2x),\cos(3x)\}$ are orthogonal over the interval $[-\pi,\pi]$.
\item Evaluate $\int_{-\pi}^{\pi}\sin(mx)\cos(nx)\,dx$.
\item Integrate $y'=\sqrt{\tan x}\sec^{4}x$.
\end{enumerate}

For each pair of integrals, determine which one is more difficult to evaluate. Explain your reasoning.

\begin{enumerate}[resume]
\item $\int \sin^{456}x\cos x\,dx$ or $\int \sin^{2}x\cos^{2}x\,dx$
\item $\int \tan^{350}x\sec^{2}x\,dx$ or $\int \tan^{350}x\sec x\,dx$
\end{enumerate}

\end{document}
